\section{Lemma 12: approximating z-rotations with H and T}
\label{sec:lemma12}

This section is the second half of the universality proof and the
mathematically deepest part of the formalization.  The goal is the paper's
Lemma \texttt{approximation-of-rz}: for every angle $\theta$ and every
$\epsilon > 0$ there is a one-qubit circuit over $\{H,T\}$ within $\epsilon$ of
$R_z(\theta)$.

The strategy is the one from \texttt{universal\_new\_gates.tex}.  Two
particular $\{H,T\}$ words generate rotations by a common \emph{irrational}
angle about two \emph{orthogonal} axes.  Orthogonality makes a three-factor
Euler decomposition valid, so $R_z(\theta)$ is exactly a product of three
rotations about those two axes; density of the integer powers of an irrational
rotation then approximates each factor.

Every node in this section is proved.  As explained in the overview, the
choice of generators is exactly what the July 2026 revision of the paper
changed: for the earlier generators the axes are not orthogonal and the
three-factor step is false.

\subsection{Axis rotations}

\begin{definition}
  \label{def:axis_rotation}
  \lean{TwoControl.Clifford.Lemma12.axisRotation}
  \leanok
  For a unit vector $n \in \mathbb{R}^3$ and $\varphi \in \mathbb{R}$, set
  $\vec{n}\cdot\vec{\sigma} = n_x\sigma_x + n_y\sigma_y + n_z\sigma_z$ and
  \[
    R(n,\varphi) \;=\; \exp\!\big(i\varphi\, \vec{n}\cdot\vec{\sigma}\big).
  \]
\end{definition}

\begin{lemma}
  \label{lem:pauli_sq}
  \lean{TwoControl.Clifford.Lemma12.pauliVec_sq_eq_one}
  \leanok
  For a unit vector $n$, $(\vec{n}\cdot\vec{\sigma})^2 = I$.
\end{lemma}

\begin{proof}
  \leanok
  Expand the square.  The diagonal terms contribute
  $(n_x^2+n_y^2+n_z^2)I = I$ and the cross terms cancel in pairs because
  distinct Pauli matrices anticommute.
\end{proof}

\begin{lemma}
  \label{lem:axis_closed_form}
  \lean{TwoControl.Clifford.Lemma12.axisRotation_closed_form}
  \leanok
  For a unit vector $n$,
  \[
    R(n,\varphi) = \cos\varphi \cdot I + i\sin\varphi \cdot
      (\vec{n}\cdot\vec{\sigma}),
  \]
  and $R(n,\varphi)$ is unitary.
\end{lemma}

\begin{proof}
  \leanok
  \uses{def:axis_rotation, lem:pauli_sq}
  By Lemma~\ref{lem:pauli_sq} the exponent squares to a scalar, so the
  exponential series splits into even and odd parts giving $\cos$ and
  $i\sin$.
\end{proof}

\begin{lemma}
  \label{lem:rz_axis}
  \lean{TwoControl.Clifford.Lemma12.rz_eq_axisRotation_standardZ}
  \leanok
  $R_z(\theta) = R(e_z, -\theta/2)$, where $e_z$ is the standard $z$ axis.
\end{lemma}

\begin{proof}
  \leanok
  \uses{lem:axis_closed_form}
  Both sides are diagonal; compare entries using
  Lemma~\ref{lem:axis_closed_form}.
\end{proof}

\begin{lemma}
  \label{lem:trace_axis}
  \lean{TwoControl.Clifford.Lemma12.trace_axisRotation}
  \leanok
  For a unit vector $n$, $\operatorname{Tr} R(n,\varphi) = 2\cos\varphi$.  More
  generally, for unit vectors $n_1, n_2$,
  \[
    \operatorname{Tr}\big(R(n_1,\varphi_1)R(n_2,\varphi_2)\big)
      = 2\big(\cos\varphi_1\cos\varphi_2
        - \sin\varphi_1\sin\varphi_2\,\langle n_1, n_2\rangle\big).
  \]
\end{lemma}

\begin{proof}
  \leanok
  \uses{lem:axis_closed_form}
  The Pauli matrices are traceless, so only the scalar part of
  Lemma~\ref{lem:axis_closed_form} survives in the first claim.  For the
  product, expand both factors and use
  $\operatorname{Tr}\big((\vec{n_1}\cdot\vec\sigma)(\vec{n_2}\cdot\vec\sigma)\big)
   = 2\langle n_1,n_2\rangle$.
\end{proof}

\begin{lemma}[Every special unitary is an axis rotation]
  \label{lem:su2_axis}
  \lean{TwoControl.Clifford.Lemma12.G1G2.su2_eq_axisRotation}
  \leanok
  If $U$ is a one-qubit unitary with $\det U = 1$, then there are a unit vector
  $n$ and an angle $\alpha \in [0,\pi]$ with $U = R(n,\alpha)$.
\end{lemma}

\begin{proof}
  \leanok
  \uses{lem:axis_closed_form}
  This is Nielsen--Chuang Exercise 4.8.  Determinant one forces $U$ to have the
  form $aI + i(b\sigma_x + c\sigma_y + d\sigma_z)$ with $a$ real and
  $a^2 + b^2 + c^2 + d^2 = 1$; take $\alpha = \arccos a$ and $n$ the
  normalization of $(b,c,d)$, matching Lemma~\ref{lem:axis_closed_form}.
\end{proof}

\subsection{The two generators}

\begin{definition}
  \label{def:g1}
  \lean{TwoControl.Clifford.Lemma12.G1G2.g1}
  \leanok
  $G_1 = e^{-3i\pi/8}\, THTHT$.
\end{definition}

\begin{definition}
  \label{def:g2}
  \lean{TwoControl.Clifford.Lemma12.G1G2.g2}
  \leanok
  $G_2 = (HT^4)\, G_1\, (HT^4)^\dagger$.
\end{definition}

\begin{definition}
  \label{def:g_words}
  \lean{TwoControl.Clifford.Lemma12.G1G2.g1Word}
  \leanok
  The $\{H,T\}$ words realizing the generators are
  $w_1 = THTHT$ and $w_2 = (HT^4)\,w_1\,(T^4H)$.  They realize $G_1, G_2$ only
  up to the scalar $e^{-3i\pi/8}$, which is harmless because the
  Hilbert--Schmidt distance is phase blind (Lemma~\ref{lem:hs_smul}).
\end{definition}

\begin{lemma}
  \label{lem:g_word_eval}
  \lean{TwoControl.Clifford.Lemma12.G1G2.g1Word_eval}
  \leanok
  $\operatorname{eval}(w_1) = e^{3i\pi/8}\, G_1$ and
  $\operatorname{eval}(w_2) = e^{3i\pi/8}\, G_2$.
\end{lemma}

\begin{proof}
  \leanok
  \uses{def:g1, def:g2, def:g_words}
  Immediate from the definitions: the word evaluates to the raw matrix product,
  and the generator is that product times the phase.
\end{proof}

\begin{lemma}
  \label{lem:g_det}
  \lean{TwoControl.Clifford.Lemma12.G1G2.g1_det}
  \leanok
  $\det G_1 = \det G_2 = 1$, and both are unitary.
\end{lemma}

\begin{proof}
  \leanok
  \uses{def:g1, def:g2}
  $\det T = e^{i\pi/4}$ and $\det H = -1$, so the raw word $THTHT$ has
  determinant $e^{3i\pi/4}$; the scalar $e^{-3i\pi/8}$ contributes its square,
  and the two cancel.  For $G_2$ conjugation preserves the determinant.
\end{proof}

\begin{lemma}[Trace of the generators]
  \label{lem:g_trace}
  \lean{TwoControl.Clifford.Lemma12.G1G2.g1_trace_sq}
  \leanok
  $(\operatorname{Tr} G_1)^2 = (\operatorname{Tr} G_2)^2 = 1 + \tfrac{\sqrt 2}{2}$.
\end{lemma}

\begin{proof}
  \leanok
  \uses{def:g1, def:g2}
  A finite computation in $\mathbb{Q}(i,\sqrt 2)$.  Note that the formalization
  pins down the \emph{square} of the trace, never the sign: the paper's
  $\operatorname{Tr} G_i = \sqrt{1 + 1/\sqrt2}$ is not needed, since every later
  step consumes only the square.  Conjugation gives
  $\operatorname{Tr} G_2 = \operatorname{Tr} G_1$.
\end{proof}

\begin{lemma}
  \label{lem:g1g2_trace}
  \lean{TwoControl.Clifford.Lemma12.G1G2.g1_g2_trace}
  \leanok
  $\operatorname{Tr}(G_1 G_2) = \dfrac{2 + \sqrt 2}{4}$.
\end{lemma}

\begin{proof}
  \leanok
  \uses{def:g1, def:g2}
  A finite computation in $\mathbb{Q}(i,\sqrt 2)$.
\end{proof}

\subsection{The rotation angle is an irrational multiple of a full turn}

\begin{lemma}
  \label{lem:lambda_irrational}
  \lean{TwoControl.Clifford.Lemma12.G1G2.irrational_div_two_pi_of_four_cos_sq}
  \leanok
  If $\alpha$ satisfies $(2\cos\alpha)^2 = 1 + \tfrac{\sqrt2}{2}$, then
  $\alpha/(2\pi)$ is irrational.
\end{lemma}

\begin{proof}
  \leanok
  Suppose $\alpha/(2\pi) = p/q$.  Then $\zeta = e^{i\alpha}$ is a root of unity
  and hence an algebraic integer, so $x = \zeta + \zeta^{-1} = 2\cos\alpha$ is
  too, and therefore so is $(x^2-1)^2$.  The hypothesis gives
  $x^2 - 1 = 1/\sqrt2$, hence $(x^2-1)^2 = 1/2$, which is not an algebraic
  integer because $\mathbb{Z}$ is integrally closed in $\mathbb{Q}$.

  This is the paper's Lemma \texttt{a-specific-lambda-is-irrational} in
  sign-free form.  The paper's constant
  $\lambda = \tfrac1\pi \arccos\!\big(\tfrac12\sqrt{1+1/\sqrt2}\big)$ is never
  named and the sign of $\cos\alpha$ is never determined; the argument consumes
  only the square.
\end{proof}

\begin{lemma}
  \label{lem:g_axis_data}
  \lean{TwoControl.Clifford.Lemma12.G1G2.g1_axis_data}
  \leanok
  There are a unit vector $n_1$ and an angle $\alpha_1$ with
  $G_1 = R(n_1,\alpha_1)$, $\cos^2\alpha_1 = \tfrac{2+\sqrt2}{8}$,
  $\sin\alpha_1 > 0$, and $\alpha_1/(2\pi)$ irrational; likewise for $G_2$ with
  the same squared cosine.
\end{lemma}

\begin{proof}
  \leanok
  \uses{lem:su2_axis, lem:g_det, lem:g_trace, lem:trace_axis,
        lem:lambda_irrational}
  Lemmas~\ref{lem:g_det} and~\ref{lem:su2_axis} put $G_i$ in the form
  $R(n_i,\alpha_i)$ with $\alpha_i \in [0,\pi]$, so $\sin\alpha_i \ge 0$.  By
  Lemma~\ref{lem:trace_axis}, $\operatorname{Tr} G_i = 2\cos\alpha_i$, and
  Lemma~\ref{lem:g_trace} pins $(2\cos\alpha_i)^2 = 1 + \sqrt2/2$, which is the
  stated squared cosine and is not $4$, so $\sin\alpha_i > 0$.
  Lemma~\ref{lem:lambda_irrational} supplies the irrationality.
\end{proof}

\begin{lemma}[The axes are orthogonal]
  \label{lem:g_axes}
  \lean{TwoControl.Clifford.Lemma12.G1G2.g_axes_data}
  \leanok
  With $G_1 = R(n_1,\alpha_1)$ and $G_2 = R(n_2,\alpha_2)$ as above,
  $\langle n_1, n_2\rangle = 0$.
\end{lemma}

\begin{proof}
  \leanok
  \uses{lem:g_axis_data, lem:g1g2_trace, lem:trace_axis}
  By Lemma~\ref{lem:trace_axis},
  $\operatorname{Tr}(G_1G_2)
   = 2(\cos\alpha_1\cos\alpha_2
       - \sin\alpha_1\sin\alpha_2\langle n_1,n_2\rangle)$, and
  Lemma~\ref{lem:g1g2_trace} evaluates the left side to $(2+\sqrt2)/4$.  The
  squared cosines from Lemma~\ref{lem:g_axis_data} force
  $\cos\alpha_1\cos\alpha_2 = \pm\tfrac{2+\sqrt2}{8}$.  In the $+$ case the
  equation reduces to $\sin\alpha_1\sin\alpha_2\langle n_1,n_2\rangle = 0$ with
  both sines positive, so $\langle n_1,n_2\rangle = 0$.  In the $-$ case the
  equation forces $\lvert\langle n_1,n_2\rangle\rvert > 1$, contradicting
  Cauchy--Schwarz for unit vectors.

  This is the paper's Lemma \texttt{a1-and-a2-anticommute}.  It is obtained
  from two trace computations rather than by exhibiting the Hermitian
  generators, and it is the fact that makes the Euler step below valid.
\end{proof}

\subsection{Euler decomposition over two orthogonal axes}

\begin{lemma}
  \label{lem:euler_expansion}
  \lean{TwoControl.Clifford.Lemma12.euler_product_expansion}
  \leanok
  For unit vectors $n_1 \perp n_2$ and angles $\alpha,\beta,\gamma$, the
  product $R(n_1,\alpha)R(n_2,\beta)R(n_1,\gamma)$ equals the special unitary
  with scalar part $\cos\beta\cos(\gamma+\alpha)$ and vector part
  \[
    \cos\beta\sin(\gamma{+}\alpha)\,n_1
      + \sin\beta\cos(\gamma{-}\alpha)\,n_2
      + \sin\beta\sin(\gamma{-}\alpha)\,(n_1\times n_2).
  \]
\end{lemma}

\begin{proof}
  \leanok
  \uses{lem:axis_closed_form}
  Multiply out the three closed forms of Lemma~\ref{lem:axis_closed_form}.
  Orthogonality makes $(n_1, n_2, n_1\times n_2)$ an orthonormal frame, in
  which the Pauli products simplify to the stated combination.
\end{proof}

\begin{lemma}[Generalized Euler decomposition]
  \label{lem:euler_orthogonal}
  \lean{TwoControl.Clifford.Lemma12.standardZ_axisRotation_orthogonal_euler}
  \leanok
  For unit vectors $n_1 \perp n_2$ and any $\varphi$, there exist
  $a, b, c$ with
  \[
    R(e_z, \varphi) = R(n_1, a)\, R(n_2, b)\, R(n_1, c).
  \]
\end{lemma}

\begin{proof}
  \leanok
  \uses{lem:euler_expansion, lem:axis_closed_form}
  This is Nielsen--Chuang Exercise 4.11 for orthogonal axes.  Expand the
  left-hand side in the orthonormal frame $(n_1,n_2,n_1\times n_2)$ and solve
  for $a, b, c$ using Lemma~\ref{lem:euler_expansion}: $b$ is recovered from
  the component along $n_1$, and $\gamma \pm \alpha$ from two argument
  extractions.  No constraint on $\varphi$ is needed --- this is exactly the
  step that fails for non-orthogonal axes.
\end{proof}

\subsection{Density of integer powers}

\begin{definition}
  \label{def:zpow_circuit}
  \lean{TwoControl.Clifford.Lemma12.zpowCircuit}
  \leanok
  For an $\{H,T\}$ word $C$ and $k \in \mathbb{Z}$, the word $C^{k}$ is
  $C$ repeated $\lvert k\rvert$ times, using the letterwise inverse word for
  negative $k$.  Inverses stay inside $\{H,T\}$ because $T^{-1} = T^7$ and
  $H^{-1} = H$; concretely $G_1^{-1}$ is realized by $T^7HT^7HT^7$.
\end{definition}

\begin{lemma}
  \label{lem:zpow_eval}
  \lean{TwoControl.Clifford.Lemma12.eval_zpowCircuit}
  \leanok
  $\operatorname{eval}(C^{k})$ is the $k$-th integer power of
  $\operatorname{eval}(C)$.
\end{lemma}

\begin{proof}
  \leanok
  \uses{def:zpow_circuit}
  Induct on $\lvert k \rvert$; the negative case uses that the letterwise
  inverse word evaluates to the inverse matrix.
\end{proof}

\begin{lemma}
  \label{lem:hs_smul}
  \lean{TwoControl.Clifford.Lemma12.hsDistance_smul_right}
  \leanok
  $d(U, zV) = d(U,V)$ whenever $\lvert z\rvert = 1$.
\end{lemma}

\begin{proof}
  \leanok
  \uses{def:hs_distance}
  Scaling by a unit scalar multiplies the trace by that scalar and so does not
  change its modulus.
\end{proof}

\begin{lemma}[Density]
  \label{lem:dense}
  \lean{TwoControl.Clifford.Lemma12.axisRotation_powers_dense}
  \leanok
  Let $n$ be a unit vector and $\theta$ an angle with $\theta/(2\pi)$
  irrational.  Then for every target angle $\alpha$ and every $\epsilon > 0$
  there is $k \in \mathbb{Z}$ with
  $d\big(R(n,\alpha), R(n,\theta)^{k}\big) < \epsilon$.
\end{lemma}

\begin{proof}
  \leanok
  \uses{lem:axis_closed_form, def:hs_distance}
  $R(n,\theta)^k = R(n,k\theta)$, so the claim is that
  $\{k\theta \bmod 2\pi\}$ comes arbitrarily close to $\alpha$.  Irrationality
  of $\theta/(2\pi)$ makes the subgroup generated by $\theta$ and $2\pi$ dense
  in $\mathbb{R}$ (Hardy--Wright; available in Mathlib as the density criterion
  for the additive subgroup generated by two reals with irrational ratio), and
  $\varphi \mapsto R(n,\varphi)$ is continuous.
\end{proof}

\begin{lemma}[Density, phased form]
  \label{lem:dense_smul}
  \lean{TwoControl.Clifford.Lemma12.axisRotation_powers_dense_smul}
  \leanok
  The same conclusion holds for $U = z\,R(n,\theta)$ with
  $\lvert z\rvert = 1$: for every $\alpha$ and $\epsilon > 0$ there is
  $k \in \mathbb{Z}$ with $d\big(R(n,\alpha), U^{k}\big) < \epsilon$.
\end{lemma}

\begin{proof}
  \leanok
  \uses{lem:dense, lem:hs_smul}
  $U^k = z^k R(n,\theta)^k$ with $\lvert z^k\rvert = 1$; apply
  Lemma~\ref{lem:hs_smul} and then Lemma~\ref{lem:dense}.  This is what lets
  the argument use the $\{H,T\}$ words directly, phase and all.
\end{proof}

\begin{lemma}
  \label{lem:hs_triple}
  \lean{TwoControl.Clifford.Lemma12.hsDistance_triple_mul_le}
  \leanok
  For unitaries, $d(ABC, A'B'C') \le d(A,A') + d(B,B') + d(C,C')$.
\end{lemma}

\begin{proof}
  \leanok
  \uses{lem:hs_mul}
  Apply Lemma~\ref{lem:hs_mul} twice.
\end{proof}

\subsection{Lemma 12}

\begin{theorem}[$R_z$ approximation by $\{H,T\}$, core form]
  \label{thm:lemma12_core}
  \lean{TwoControl.Clifford.Lemma12.G1G2.HT_rz_dense_g1g2}
  \leanok
  For every $\theta$ and every $\epsilon > 0$ there is an $\{H,T\}$ word $C$
  with $d\big(R_z(\theta), \operatorname{eval}(C)\big) < \epsilon$.
\end{theorem}

\begin{proof}
  \leanok
  \uses{lem:g_axes, lem:euler_orthogonal, lem:rz_axis, lem:dense_smul,
        lem:hs_triple, lem:g_word_eval, lem:zpow_eval}
  Take the axis data of Lemma~\ref{lem:g_axes}: unit axes $n_1 \perp n_2$,
  irrational angles $\alpha_1, \alpha_2$, and $G_i = R(n_i,\alpha_i)$.

  By Lemma~\ref{lem:rz_axis}, $R_z(\theta) = R(e_z, -\theta/2)$, and
  Lemma~\ref{lem:euler_orthogonal} produces angles $a,b,c$ with
  \[
    R_z(\theta) = R(n_1,a)\,R(n_2,b)\,R(n_1,c),
  \]
  an \emph{exact} identity.  Lemma~\ref{lem:g_word_eval} shows the words $w_1,
  w_2$ evaluate to $G_1, G_2$ up to a unit scalar, so
  Lemma~\ref{lem:dense_smul} gives integers $k_1,k_2,k_3$ with each of
  $w_1^{k_1}, w_2^{k_2}, w_1^{k_3}$ within $\epsilon/3$ of the corresponding
  Euler factor, using Lemma~\ref{lem:zpow_eval} to identify the evaluation of a
  power word with the matrix power.  Concatenating the three words and applying
  Lemma~\ref{lem:hs_triple} bounds the total error by $\epsilon$.
\end{proof}

\begin{theorem}[Lemma 12]
  \label{thm:lemma12}
  \lean{TwoControl.Clifford.Lemma12.lemma12_rz_approximation_by_ht}
  \leanok
  For every $\theta$ and every $\epsilon > 0$ there is a list of $\{H,T\}$
  primitives whose circuit matrix is within $\epsilon$ of $R_z(\theta)$ in
  Hilbert--Schmidt distance.
\end{theorem}

\begin{proof}
  \leanok
  \uses{thm:lemma12_core, def:ht_primitive}
  Restatement of Theorem~\ref{thm:lemma12_core} in the circuit vocabulary of
  Definition~\ref{def:ht_primitive}, which is the interface the $n$-qubit
  assembly consumes.
\end{proof}
